\documentclass[UTF8,12pt]{ctexart}
\usepackage[a4paper,margin=2.3cm]{geometry}
\usepackage{amsmath,amssymb}
\usepackage{array}
\usepackage{enumitem}
\usepackage{fancyhdr}
\usepackage{lastpage}
\usepackage[colorlinks,
linkcolor=black,
anchorcolor=black,
citecolor=black
]{hyperref}

\setlength{\parindent}{2em}
\setlength{\parskip}{0.35em}
\linespread{1.15}
\pagestyle{fancy}
\fancyhf{}
\cfoot{第 \thepage\ 页，共 \pageref{LastPage} 页}
\renewcommand{\headrulewidth}{0pt}

\newcommand{\blank}[1][4em]{\underline{\hspace{#1}}}
\newcommand{\examsection}[1]{\par\addvspace{1em}\noindent{\large\bfseries #1}\par\addvspace{0.4em}}
\newcommand{\choice}[4]{%
  \begin{tabular}{@{}p{0.48\linewidth}p{0.48\linewidth}@{}}
  (A) #1 & (B) #2\\[0.4em]
  (C) #3 & (D) #4
  \end{tabular}
}

\begin{document}

\begin{center}
  {\LARGE\bfseries 《线性代数》期中考试试题}\\[0.5em]
  {\large 2025--2026 学年 第二学期 / 福建师范大学数学与统计学院 }\\[0.3em]
  排版：\href{https://github.com/Xuuyuan}{@许愿} -- \href{https://wefjnu.nekoark.com}{WeFJNU}
\end{center}

\noindent 说明：试卷中的矩阵默认都是实矩阵，$E$ 为单位矩阵。

\examsection{一、单项选择题（每小题 3 分，共 18 分）}
\begin{enumerate}[label={\arabic*、},leftmargin=2em,itemsep=0.8em]
\item 设 $A,B,C$ 都是 $n$ 阶方阵，则下列结论正确的是（\quad）

\choice{$(A+B)^{-1}=A^{-1}+B^{-1}$；}
{若 $AC=BC$ 且 $C\ne O$，则 $A=B$；}
{$|ABC|=|CBA|$；}
{$A^2-B^2=(A-B)(A+B)$。}

\item 设
\[
D_1=\begin{vmatrix}
a_{11}&a_{12}&a_{13}\\
a_{21}&a_{22}&a_{23}\\
a_{31}&a_{32}&a_{33}
\end{vmatrix},\quad
D_2=\begin{vmatrix}
2a_{11}&a_{12}-3a_{13}&a_{13}\\
2a_{21}&a_{22}-3a_{23}&a_{23}\\
2a_{31}&a_{32}-3a_{33}&a_{33}
\end{vmatrix},
\]
则以下结论正确的是（\quad）

\choice{$D_2=-3D_1$；}{$D_2=2D_1$；}{$D_2=-6D_1$；}{$D_2=5D_1$。}

\item 若 $4\times 5$ 矩阵 $A$ 的秩为 3，以下说法\textbf{未必}正确的是（\quad）

\choice{$A$ 仅有一个三阶不为零的子式；}
{$A$ 至少有一个二阶子式不等于零；}
{$A$ 的所有 4 阶子式都等于 0；}
{$A$ 和 $\begin{pmatrix}O&O\\O&E_3\end{pmatrix}$ 等价。}

\item 设矩阵 $A=(\alpha_1,\alpha_2,\alpha_3)$ 经过初等行变换可化为
$\begin{pmatrix}1&1&3\\0&1&1\end{pmatrix}$，则下列说法\textbf{错误}的是（\quad）

\choice{齐次线性方程组 $AX=0$ 有非零解；}
{$R(A)=2$；}
{$A$ 的等价标准形为 $\begin{pmatrix}1&0&0\\0&1&0\end{pmatrix}$；}
{齐次线性方程组 $A^{\mathrm T}X=0$ 有非零解。}

\item 设 $n$ 阶方阵 $A$ 满足 $R(A)<n$，则下列说法\textbf{错误}的是（\quad）

\begin{enumerate}[label=(\Alph*),leftmargin=2em,itemsep=0.2em]
\item $A$ 中可能有一行元素全为 0；
\item $A$ 中可能有两列元素对应成比例；
\item 齐次线性方程组 $AX=0$ 有非零解；
\item 非齐次线性方程组 $AX=\beta$ 可能有唯一解。
\end{enumerate}

\newpage
\item 设矩阵 $A=(2)$，$B=\begin{pmatrix}1&1\\1&4\end{pmatrix}$，则分块矩阵
$\begin{pmatrix}O&A\\B&O\end{pmatrix}$ 的伴随矩阵
$\begin{pmatrix}O&A\\B&O\end{pmatrix}^{*}$ 为（\quad）

\begin{tabular}{@{}p{0.48\linewidth}p{0.48\linewidth}@{}}
(A) $\begin{pmatrix}0&8&-2\\0&-2&2\\3&0&0\end{pmatrix}$ &
(B) $\begin{pmatrix}0&0&3\\8&-2&0\\-2&2&0\end{pmatrix}$ \\[1.6em]
(C) $\begin{pmatrix}0&0&\frac12\\\frac43&-\frac13&0\\-\frac13&\frac13&0\end{pmatrix}$ &
(D) $\begin{pmatrix}0&\frac43&-\frac13\\0&-\frac13&\frac13\\\frac12&0&0\end{pmatrix}$
\end{tabular}
\end{enumerate}

\examsection{二、填空题（每小题 3 分，共 18 分）}
\begin{enumerate}[label={\arabic*、},leftmargin=2em,itemsep=0.8em]
\item 已知 2 阶方阵 $X$ 满足 $\begin{pmatrix}1&0\\-1&1\end{pmatrix}X\begin{pmatrix}0&1\\1&0\end{pmatrix}^{2026}=\begin{pmatrix}3&1\\0&-1\end{pmatrix}$，则 $X=$ \blank[8em]。

\item 设行列式
\[
|A|=\begin{vmatrix}
1&0&1&0\\
2&0&4&1\\
1&3&4&2\\
4&3&3&-1
\end{vmatrix},
\]
$M_{ij}$ 和 $A_{ij}$ $(i,j=1,2,3,4)$ 分别为 $|A|$ 的 $(i,j)$ 元的余子式和代数余子式，则
$2A_{31}+3A_{32}-2M_{34}=$ \blank[8em]。

\item 设 $A$ 为 3 阶方阵且满足 $A^2-2A-6E=0$，则 $(A-3E)^{-1}=$ \blank[8em]。

\item 设 $A$ 为 4 阶方阵且 $|A|=3$，则 $\left|A+AE(1,3(2))\right|=$ \blank[8em]。

\item 设矩阵
\[
A=\begin{pmatrix}1&1&1\\0&2&1\\0&0&3\end{pmatrix},\quad
B=\begin{pmatrix}1&1&1\\2&3&4\\3&3&3\end{pmatrix},\quad
C=\begin{pmatrix}0&0&1\\0&2&1\\3&1&1\end{pmatrix},
\]
则 $R(ABC)=$ \blank[8em]。

\item 已知齐次线性方程组 $\begin{cases}x_1+x_2+x_3=0,\\x_1+2x_2+ax_3=0,\\x_1+4x_2+a^2x_3=0\end{cases}$ 有非零解，则 $a=$ \blank[8em]。
\end{enumerate}

\newpage
\examsection{三、（12 分）}
设矩阵
\[
A=\begin{pmatrix}2&3&0\\1&2&0\\0&0&-1\end{pmatrix},\quad
B=\begin{pmatrix}1&1&0\\2&2&0\\0&0&1\end{pmatrix},
\]
求 $\left|AB^{\mathrm T}-2E\right|$。

\examsection{四、（20 分）}
已知 $\alpha=(1,0,2)^{\mathrm T}$，$\beta=(0,1,0)^{\mathrm T}$，$A=E+\alpha\beta^{\mathrm T}$。
\begin{enumerate}[label=(\arabic*),leftmargin=2em,itemsep=0.6em]
\item 证明 $A$ 可逆并求 $A^{-1}$；
\item 将 $A$ 的伴随矩阵 $A^{*}$ 按列分块成 $A^{*}=(A_1,A_2,A_3)$，求 $A(A_1+2A_2+3A_3)$；
\item 若 $f(x)=5(x-1)^3-3(x-1)^2+2(x-1)-4$，求 $f(A)$。
\end{enumerate}

\examsection{五、（12 分）}
已知
\[
A=\begin{pmatrix}1&1&-1\\-1&1&1\\1&-1&1\end{pmatrix},
\]
$A^{*}$ 是 $A$ 的伴随矩阵，矩阵 $X$ 满足 $A^{*}X=A^{-1}+2X$，求矩阵 $X$。

\examsection{六、（12 分）}
用初等行变换求解非齐次线性方程组 $\begin{cases}x_1-x_2+2x_4=-1,\\2x_1-2x_2-x_3+2x_4=-2,\\x_1-2x_2+x_3+x_4=2,\\3x_1-4x_2+x_3+5x_4=0.\end{cases}$

\examsection{七、（8 分）}
如果 $n$ 阶方阵 $A$ 满足 $A^2=A$，则称 $A$ 是一个\textbf{幂等矩阵}。已知 $n$ 阶方阵 $A,B$ 都是幂等矩阵，证明：$A+B$ 是幂等矩阵的充分必要条件是 $AB=BA=O$。

\end{document}
